\documentclass[12pt]{article}

\usepackage{amsmath,amssymb,amsthm}
\usepackage{geometry}
\usepackage{booktabs}
\usepackage{array}
\usepackage{natbib}
\usepackage{hyperref}

\geometry{margin=1in}

\newtheorem{proposition}{Proposition}
\newtheorem{lemma}{Lemma}
\newtheorem{remark}{Remark}

\setlength{\parskip}{0.5em}
\setlength{\parindent}{0pt}

\title{Sequential Monitoring for Structural Breaks in Heterogeneous Panels:\\
Design, Derivations, and Computational Details}

\author{}
\date{}

\begin{document}

\maketitle

\begin{abstract}
\noindent
This note describes a Monte Carlo study of a real-time structural break
monitoring procedure for heterogeneous panels. The procedure works in three
steps: (1) each series is pre-filtered by an autoregression estimated on a
calibration sample; (2) the filtered residuals are aggregated across series
via $N(N-1)$ separate (unpooled) pairwise regressions, one OLS fit per
ordered pair of series; (3) the cumulative sum of squares (CUSQ) of
the aggregate residual is compared period by period against a boundary
derived from the Brownian bridge, following \citet{WangHsiao2013}. The note
gives the exact data generating processes, a closed-form derivation showing
that the $N(N-1)$-pair aggregation, even though every pair has its own
fitted slope, can be evaluated exactly as a weighted cross-sectional average
without ever materializing the $N(N-1)$ residual series, the two-scenario
Monte Carlo design (size under no break, power and
detection delay under a known break), and two computational improvements
--- FFT-based fractional differencing and vectorized AR residuals --- that
make the full $1.9\times10^6$-replication design feasible without changing
any statistic. A companion design adds four targeted power experiments ---
common-factor and subgroup-specific volatility breaks, and common-factor
and subgroup-specific loading breaks --- built on the same AR-filtering,
pairwise-aggregation, and CUSQ machinery.
\end{abstract}

\section{Setup}
\label{sec:setup}

We observe a panel of $N$ time series $\{Y_{i,t}\}$ for $i=1,\dots,N$ and
$t=1,\dots,T$. The first $M$ periods form the calibration sample, assumed to
be free of breaks. From period $M+1$ onward the monitor checks, one period at
a time, whether the panel has shifted.

The null and alternative hypotheses are
\[
H_0:\ \text{no break for any } t > M,
\]
\[
H_1:\ \text{all } N \text{ series shift by } \delta \neq 0
      \text{ at some time } \tau > M.
\]

Sections~\ref{sec:dgp}--\ref{sec:cusum} describe the four steps of the
procedure: data generation, AR filtering, cross-sectional aggregation, and
the CUSQ test. Section~\ref{sec:mc} describes the Monte Carlo design,
Section~\ref{sec:targeted} describes four additional targeted power
experiments built on the same monitoring procedure, and
Section~\ref{sec:computation} documents the computational improvements.

\section{Data Generating Processes}
\label{sec:dgp}

\subsection{Panel series}

Each series follows an ARFIMA$(p_i,d_i,q_i)$ model,
\[
\Phi_i(L)(1-L)^{d_i}Y_{i,t} = \Theta_i(L)\,\varepsilon_{i,t},
\qquad \varepsilon_{i,t} \sim \text{i.i.d.}\; N(0,1),
\]
where $\Phi_i(L)=1-\phi_i L$ and $\Theta_i(L)=1+\theta_i L$ are at most
first-order polynomials and $(1-L)^{d_i}$ is the fractional differencing
operator. Two panel sizes are used.

\paragraph{$N=10$.} Five pairs of series, each pair sharing the same
parameters (Table~\ref{tab:N10}).

\begin{table}[h]
\centering
\caption{DGP for $N=10$.}
\label{tab:N10}
\begin{tabular}{ccccc}
\toprule
Series & Model & $\phi_i$ & $d_i$ & $\theta_i$ \\
\midrule
$1,2$  & ARMA$(1,1)$     & $0.8$ & $0$   & $0.5$ \\
$3,4$  & ARMA$(1,1)$     & $0.5$ & $0$   & $0.3$ \\
$5,6$  & ARFIMA$(0,d,0)$ & $0$   & $0.4$ & $0$   \\
$7,8$  & ARFIMA$(1,d,1)$ & $0.1$ & $0.3$ & $0.8$ \\
$9,10$ & ARFIMA$(1,d,0)$ & $0.7$ & $0.4$ & $0$   \\
\bottomrule
\end{tabular}
\end{table}

\paragraph{$N=30$.} A random-coefficient panel divided into four blocks
(Table~\ref{tab:N30}). The coefficients are drawn independently for each
series and for each Monte Carlo replication. This means the simulation
averages over both the randomness in the coefficients and the randomness in
the innovations, giving a measure of average performance across the whole
class of DGPs described in the table rather than performance at one fixed
parameter vector.

\begin{table}[h]
\centering
\caption{DGP for $N=30$ (random coefficients, redrawn each replication).}
\label{tab:N30}
\begin{tabular}{ccccc}
\toprule
Series & Model & $\phi_i$ & $d_i$ & $\theta_i$ \\
\midrule
$1$--$15$  & ARMA$(1,1)$     & $U(0,\,0.5)$   & $0$            & $0.6$ \\
$16$--$20$ & ARMA$(1,1)$     & $U(0.7,\,0.9)$ & $0$            & $0.6$ \\
$21$--$25$ & ARFIMA$(0,d,0)$ & $0$            & $U(0,\,0.5)$   & $0$   \\
$26$--$30$ & ARFIMA$(1,d,1)$ & $U(0,\,0.4)$   & $U(0,\,0.5)$   & $U(0,\,0.3)$ \\
\bottomrule
\end{tabular}
\end{table}

\subsection{Common factors}

Four common factors are generated independently of the panel series
(Table~\ref{tab:factors}). They cover i.i.d.\ noise, a near-unit-root AR(1),
and two long-memory processes.

\begin{table}[h]
\centering
\caption{Common factor DGPs.}
\label{tab:factors}
\begin{tabular}{cccc}
\toprule
Factor & Model & $\phi$ & $d$ \\
\midrule
1 & i.i.d.\ $N(0,1)$  & $0$    & $0$    \\
2 & AR$(1)$           & $0.99$ & $0$    \\
3 & ARFIMA$(0,d,0)$   & $0$    & $0.20$ \\
4 & ARFIMA$(0,d,0)$   & $0$    & $0.45$ \\
\bottomrule
\end{tabular}
\end{table}

\subsection{Observed data and break injection}

The observed series add a common factor with loading $\lambda = 0.5$:
\begin{equation}
Y^{*}_{i,t} = Y_{i,t} + 0.5\,f_t, \qquad i=1,\dots,N.
\label{eq:factorload}
\end{equation}
Under $H_1$, a level shift of size $\delta$ is added at time $\tau$:
\begin{equation}
Y^{**}_{i,t} = Y^{*}_{i,t} + \delta \cdot \mathbf{1}\{t \ge \tau\},
\qquad i=1,\dots,N.
\label{eq:break}
\end{equation}
The break time $\tau$ is specified in Section~\ref{subsec:scenarioB}.

\section{AR Pre-Filtering}
\label{sec:arfilter}

For each series $i$, we fit an AR($p$) model on the calibration sample
$t=1,\dots,M$, where
\begin{equation}
p = \max\!\left(3,\; 3\left\lfloor\frac{M-1}{100}\right\rfloor\right).
\label{eq:p}
\end{equation}
The OLS estimate from the calibration sample is
\begin{equation}
\hat\beta_{AR} = \arg\min_{b\in\mathbb{R}^p}
\sum_{t=p+1}^{M}
\Bigl(Y_{i,t} - b'(Y_{i,t-1},\dots,Y_{i,t-p})'\Bigr)^2.
\label{eq:arcal}
\end{equation}
This estimate is fixed after period $M$ and never updated. For all
$t = p+1,\dots,T$, the one-step-ahead residual is
\begin{equation}
\hat e_{i,t} = Y_{i,t} - \hat\beta_{AR}'(Y_{i,t-1},\dots,Y_{i,t-p})'.
\label{eq:residual}
\end{equation}
The first seven observations are discarded as a burn-in. Fixing the AR
coefficients after calibration is the key sequential monitoring requirement:
parameters are estimated from clean data and then held constant, so the
monitoring period residuals are not used to update the model.

\section{Cross-Sectional Aggregation}
\label{sec:pool}

\subsection{Unpooled pairwise regression}

Let $\hat e_{i,t}$ denote the AR residuals for $i=1,\dots,N$ and
$t=1,\dots,T$. For every ordered pair $(i,j)$ with $i\ne j$ (there are
$N(N-1)$ such pairs), we fit a \emph{separate} calibration-period
regression
\begin{equation}
\hat e_{i,t} = \alpha_{ij} + \beta_{ij}\,\hat e_{j,t} + u_{i,j,t},
\qquad t=1,\dots,M.
\label{eq:pairreg}
\end{equation}
Unlike a pooled specification, the $N(N-1)$ regressions are \emph{not}
stacked into a single sample: each pair $(i,j)$ has its own OLS estimate
$(\hat\alpha_{ij},\hat\beta_{ij})$, computed using only that pair's $M$
calibration observations. As with the AR coefficients of
Section~\ref{sec:arfilter}, $(\hat\alpha_{ij},\hat\beta_{ij})$ are estimated
once from the calibration sample and then held fixed for every
$t=1,\dots,T$. For all $t$, the residual from pair $(i,j)$'s regression is
\begin{equation}
U_{i,j,t} = \hat e_{i,t} - \hat\alpha_{ij} - \hat\beta_{ij}\,\hat e_{j,t}.
\label{eq:Uijt}
\end{equation}
The aggregate signal used in the CUSQ test is, as before, the average of
$U_{i,j,t}$ over all ordered pairs:
\begin{equation}
\bar e_t \equiv
\frac{1}{N(N-1)}\sum_{\substack{i,j=1\\i\ne j}}^{N} U_{i,j,t},
\qquad t=1,\dots,T.
\label{eq:ebar_def}
\end{equation}
As stated, computing \eqref{eq:ebar_def} requires fitting all $N(N-1)$
pairs and forming all $N(N-1)$ residual sequences --- an $O(N^2M)$
estimation step followed by an $O(N^2T)$ evaluation step. The next two
propositions show that every $(\hat\alpha_{ij},\hat\beta_{ij})$, and hence
$\bar e_t$, can be obtained in closed form at a total cost of $O(N^2M)$
(estimation) plus $O(NT)$ (evaluation), without ever materializing the
$N(N-1)$ residual series and without any approximation.

\subsection{Closed form for $(\hat\alpha_{ij},\hat\beta_{ij})$}

Define the per-series calibration mean
\begin{equation}
\mu_i \equiv \frac{1}{M}\sum_{t=1}^{M}\hat e_{i,t}, \qquad i=1,\dots,N,
\label{eq:mu}
\end{equation}
and the demeaned residual, held fixed and applied to \emph{every} $t$ (not
only the calibration sample),
\begin{equation}
f_{i,t} \equiv \hat e_{i,t} - \mu_i, \qquad t=1,\dots,T.
\label{eq:fit}
\end{equation}
Let $d_{i,t} \equiv f_{i,t}$ for $t \le M$ denote the calibration-period
values, and define the $N\times N$ Gram matrix of calibration-period
centered residuals
\begin{equation}
S_{ij} \equiv \sum_{t=1}^{M} d_{i,t}\,d_{j,t}, \qquad i,j=1,\dots,N.
\label{eq:Sij}
\end{equation}

\begin{proposition}
\label{prop:ab}
The single-pair OLS estimates from \eqref{eq:pairreg} are, for $i\ne j$,
\begin{equation}
\hat\beta_{ij} = \frac{S_{ij}}{S_{jj}}, \qquad
\hat\alpha_{ij} = \mu_i - \hat\beta_{ij}\,\mu_j.
\label{eq:abclosed}
\end{equation}
\end{proposition}

\begin{proof}
For a fixed pair $(i,j)$, \eqref{eq:pairreg} is an ordinary simple linear
regression of $\hat e_{i,t}$ on $\hat e_{j,t}$ over $t=1,\dots,M$, with no
pooling across pairs. Its OLS slope is the ratio of the sample covariance
of $(\hat e_{i,t},\hat e_{j,t})$ to the sample variance of $\hat e_{j,t}$:
$\sum_t d_{i,t}d_{j,t}\big/\sum_t d_{j,t}^2 = S_{ij}/S_{jj}$. The intercept
follows from the standard simple-regression identity
$\hat\alpha_{ij} = \mu_i - \hat\beta_{ij}\mu_j$.
\end{proof}

\subsection{Closed form for $\bar e_t$}

Define the cross-sectional sum of the centered residuals,
$F_t \equiv \sum_{i=1}^N f_{i,t}$, and, for each $j$, the column sum of
the fixed pairwise slopes excluding the diagonal,
\begin{equation}
c_j \equiv \sum_{i\ne j} \hat\beta_{ij}, \qquad j=1,\dots,N.
\label{eq:cj}
\end{equation}

\begin{proposition}
\label{prop:ebar}
For every $t=1,\dots,T$,
\begin{equation}
\bar e_t = \frac{(N-1)F_t - \displaystyle\sum_{j=1}^N c_j f_{j,t}}{N(N-1)}.
\label{eq:ebarclosed}
\end{equation}
\end{proposition}

\begin{proof}
By \eqref{eq:abclosed}, $\hat\alpha_{ij}+\hat\beta_{ij}\mu_j=\mu_i$, so
$U_{i,j,t}=\hat e_{i,t}-\hat\alpha_{ij}-\hat\beta_{ij}\hat e_{j,t}
=f_{i,t}-\hat\beta_{ij}f_{j,t}$. Expanding \eqref{eq:ebar_def}:
\[
\bar e_t = \frac{1}{N(N-1)}\sum_{i\ne j}\bigl(f_{i,t}-\hat\beta_{ij}f_{j,t}\bigr)
= \frac{1}{N(N-1)}\left[\sum_{i\ne j} f_{i,t}
  - \sum_{i\ne j}\hat\beta_{ij}f_{j,t}\right].
\]
In the first sum, $f_{i,t}$ is added once for every $j\ne i$, i.e.\ $N-1$
times for each $i$, so $\sum_{i\ne j}f_{i,t}=(N-1)F_t$. In the second sum,
regrouping by $j$ gives $\sum_{i\ne j}\hat\beta_{ij}f_{j,t}
=\sum_j f_{j,t}\sum_{i\ne j}\hat\beta_{ij}=\sum_j c_j f_{j,t}$. Combining
the two terms gives \eqref{eq:ebarclosed}.
\end{proof}

\begin{remark}
Unlike under pooling, the $N(N-1)$ pairwise regressions here do
\emph{not} share a common slope, so $\bar e_t$ is no longer a simple
affine map of the plain cross-sectional mean: each series $j$ contributes
through $f_{j,t}$ with its own pair-specific weight $c_j$.
\eqref{eq:ebarclosed} is nonetheless exact and never requires forming the
$N(N-1)$ residual series. The cost is $O(N^2M)$ to form $S$ (and hence
every $\hat\beta_{ij}$ and $c_j$) once from the calibration sample, plus
$O(NT)$ to evaluate $\bar e_t$ for the full sample --- versus $O(N^2T)$
for a direct, pair-by-pair evaluation. For $N=30$ this still replaces an
$N(N-1)=870$-pair construction with one $30\times30$ Gram matrix and a
single weighted cross-sectional average.
\end{remark}

\section{The CUSQ Statistic and Monitoring Boundary}
\label{sec:cusum}

\subsection{Statistic}

Let $C_l = \sum_{t=1}^{l}\bar e_t^2$ be the cumulative sum of squares of
$\bar e_t$. Following \citet{WangHsiao2013} (p.~5), the centered CUSQ for
$n \ge m$ is
\begin{equation}
D_n \equiv \frac{C_n}{C_m} - \frac{n-2}{m-2},
\label{eq:Dn}
\end{equation}
and the detector is
\begin{equation}
T_n \equiv \sqrt{\frac{m}{2}}\,D_n
     = \sqrt{\frac{m}{2}}\left(\frac{C_n}{C_m} - \frac{n-2}{m-2}\right).
\label{eq:Tn}
\end{equation}
This is equivalent to $T_n = m^{-1/2}\hat{S}_n$ in the notation of
\citet{WangHsiao2013}, where $\hat{S}_n = 2^{-1/2}m\hat{D}_n$.

\subsection{Limiting distribution}

Under $H_0$, \citet{WangHsiao2013} (Theorem~1) show that
\begin{equation}
T_n = m^{-1/2}\hat{S}_{[m\tau]} \;\Rightarrow\; W^0(\tau),
\qquad \tau = \frac{n}{m} \ge 1,
\label{eq:fclt}
\end{equation}
where $W^0(\tau) = W(\tau) - \tau W(1)$ is a Brownian bridge and $W(\tau)$
is a standard Wiener process.

\subsection{Monitoring boundary}

Following \citet{ChuStinchcombeWhite1996}, the crossing probability of
$W^0(\tau)$ satisfies \citep[eq.~(5)]{WangHsiao2013}:
\begin{equation}
P\bigl\{|W^0(\tau)| \ge (\tau-1)^{-1}[\tau(a^2+\ln\tau)]^{1/2}
\ \text{for some}\ \tau \ge 1\bigr\}
= 2[1-\Phi(a)] + a\phi(a),
\label{eq:crossingprob}
\end{equation}
where $\Phi$ and $\phi$ are the standard-normal CDF and PDF. Setting the
right-hand side equal to $\alpha$ and solving numerically gives
\begin{equation}
a^2_{10\%} = 6.2514, \qquad
a^2_{5\%}  = 7.8147, \qquad
a^2_{1\%}  = 11.3449.
\label{eq:a2}
\end{equation}
\citet{WangHsiao2013} state $a^2_{5\%} = 7.78$ (rounded to two decimal
places). Discretising $\lambda$ as $n/m$ yields the boundary function
\citep[p.~6]{WangHsiao2013}:
\begin{equation}
g\!\left(\frac{n}{m};\,a^2_\alpha\right)
= \frac{n-m}{m}
  \left[\frac{n}{n-m}
        \left(a^2_\alpha + \ln\frac{n}{n-m}\right)\right]^{1/2}.
\label{eq:boundary}
\end{equation}

\section{Sequential Decision Rule}
\label{sec:stopping}

At each period $n > m$, compute $T_n$ from \eqref{eq:Tn} and compare it
to the boundary $g(n/m;\,a^2_\alpha)$ from \eqref{eq:boundary}. The
procedure declares a break at the first period where the statistic crosses
the boundary:
\begin{equation}
\hat\tau_\alpha \equiv
\inf\bigl\{n > m : |T_n| \ge g\!\left(n/m;\,a^2_\alpha\right)\bigr\},
\label{eq:stoptime}
\end{equation}
with $\hat\tau_\alpha = +\infty$ if no crossing occurs before $T$. The AR
coefficients $\hat\beta_{AR}$, the pairwise regression coefficients
$(\hat\alpha_{ij},\hat\beta_{ij})$ for every $i\ne j$, and the per-series
means $\mu_i$ are all fixed from the calibration period and never updated
after $t=m$.

\section{Monte Carlo Design}
\label{sec:mc}

The total sample length used in each replication is
\begin{equation}
T_{\text{total}} = 10M + 10,
\label{eq:Ttotal}
\end{equation}
which provides a monitoring window of $9M$ periods beyond the calibration
sample. Two scenarios are run.

\subsection{Scenario A: size ($\delta=0$)}
\label{subsec:scenarioA}

No break is injected. Because the data do not depend on $k$, each
simulated path yields one stopping time $\hat\tau_\alpha^{(r)}$, and
this is scored against six window lengths
$k\in\mathcal{K}=\{0.2,0.3,0.5,1,2,3\}$ at once:
\begin{equation}
A_\alpha^{(r)}(k) \equiv
\mathbf{1}\bigl\{\hat\tau_\alpha^{(r)} \le M + kM\bigr\}.
\label{eq:Aindicator}
\end{equation}
The empirical size is
\begin{equation}
\widehat{\mathrm{Size}}_\alpha(N,M,\mathrm{factor},k)
= \frac{1}{R}\sum_{r=1}^{R} A_\alpha^{(r)}(k),
\label{eq:size}
\end{equation}
which estimates the false alarm probability within the first $kM$ monitoring
periods. A well-calibrated test should give
$\widehat{\mathrm{Size}}_\alpha \approx \alpha$.

\subsection{Scenario B: power ($\delta\in\{1,2,3\}$)}
\label{subsec:scenarioB}

A break of size $\delta$ is injected at
\begin{equation}
\tau(k) = M + \lfloor kM \rfloor, \qquad k\in\mathcal{K}.
\label{eq:tauk}
\end{equation}
Here $k$ is the break location, not a scoring window. Each $(\delta,k)$
combination requires its own set of replications because the data depend on
$(\delta,k)$. The stopping time $\hat\tau_\alpha^{(r)}$ is compared to the
true break date $\tau(k)$ and classified as one of three outcomes:
\begin{align}
P_\alpha^{(r)} &\equiv \mathbf{1}\{\hat\tau_\alpha^{(r)} \le \tau(k)\}
  && \text{(alarm before the break)}, \label{eq:premature} \\
D_\alpha^{(r)} &\equiv \mathbf{1}\{\tau(k) < \hat\tau_\alpha^{(r)} < \infty\}
  && \text{(alarm after the break)}, \label{eq:detect} \\
B_\alpha^{(r)} &\equiv P_\alpha^{(r)} + D_\alpha^{(r)}
  = \mathbf{1}\{\hat\tau_\alpha^{(r)} < \infty\}
  && \text{(any alarm).} \label{eq:anyalarm}
\end{align}
The Monte Carlo averages are
\begin{equation}
\widehat{P}_\alpha = \frac{1}{R}\sum_r P_\alpha^{(r)}, \qquad
\widehat{D}_\alpha = \frac{1}{R}\sum_r D_\alpha^{(r)}, \qquad
\widehat{B}_\alpha = \widehat{P}_\alpha + \widehat{D}_\alpha.
\label{eq:PDhats}
\end{equation}
$\widehat{D}_\alpha$ is the power of the test --- the probability of
detecting the break after it occurs. $\widehat{P}_\alpha$ is the
premature alarm rate.

For replications where $D_\alpha^{(r)}=1$, the detection delay is
\begin{equation}
\mathrm{Delay}_\alpha^{(r)} \equiv \hat\tau_\alpha^{(r)} - \tau(k) > 0.
\label{eq:delay}
\end{equation}
Summary statistics (mean, standard deviation, median, 25th and 75th
percentiles) are computed over the subset $\{r: D_\alpha^{(r)}=1\}$.
Premature alarms ($P_\alpha^{(r)}=1$) are excluded from this calculation
because they have a negative delay and measure a different phenomenon.

\subsection{Full design}

Table~\ref{tab:design} lists all parameter combinations.

\begin{table}[h]
\centering
\caption{Monte Carlo design.}
\label{tab:design}
\begin{tabular}{lll}
\toprule
Parameter & Scenario A (size) & Scenario B (power) \\
\midrule
$N$ & $\{10,30\}$ & $\{10,30\}$ \\
$\delta$ & $0$ & $\{1,2,3\}$ \\
$k$ & scoring window, $\{0.2,0.3,0.5,1,2,3\}$ & break location, $\{0.2,0.3,0.5,1,2,3\}$ \\
$M$ & $\{50,80,100,200,300\}$ & $\{50,80,100,200,300\}$ \\
Factor & $\{1,2,3,4\}$ & $\{1,2,3,4\}$ \\
$R$ & $2500$ & $2500$ \\
\midrule
Combinations & $2\times5\times4 = 40$ & $2\times3\times6\times5\times4 = 720$ \\
Total replications & $10^5$ & $1.8\times10^6$ \\
\bottomrule
\end{tabular}
\end{table}

For $N=30$, the random coefficients in Table~\ref{tab:N30} are redrawn on
every replication, so the Monte Carlo averages integrate over both the
coefficient distribution and the innovation sequences.

\section{Targeted Power Experiments}
\label{sec:targeted}

The design of Section~\ref{sec:mc} injects a level shift $\delta$ directly
into every series. A companion design instead breaks the panel
\emph{indirectly}, through the common factor that loads onto the series,
and asks whether the monitor still has power when the break works through
the factor structure rather than through a direct level shift. Four
experiments are run, crossing two break mechanisms (a break in the
factor's \emph{volatility} versus a break in its \emph{loading}) with two
factor structures (a single factor shared by the whole panel versus
factor-specific subgroups). All four reuse, unchanged, the AR pre-filtering
of Section~\ref{sec:arfilter}, the pairwise cross-sectional aggregation
\eqref{eq:ebarclosed}, and the CUSQ statistic and boundary of
Section~\ref{sec:cusum}: only the data-generating process and the reported
outcomes differ. The AR coefficients $\hat\beta_{AR}$, the pairwise
coefficients $(\hat\alpha_{ij},\hat\beta_{ij})$, and the means $\mu_i$ are
again estimated once on $t=1,\dots,M$ and held fixed thereafter. In all
four experiments, $k$ no longer plays the dual role of Section~\ref{sec:mc}
(scoring window in Scenario A, break location in Scenario B): here $k$ is
only a break location, restricted to $k\in\{0.2,0.3,0.5\}$, and the break
time is $\tau^*(k) = M + \lfloor kM\rfloor$ as in \eqref{eq:tauk}. Outcomes
are reported with the alarm/premature/detect/delay logic of
\eqref{eq:premature}--\eqref{eq:delay}, with $\tau(k)$ replaced by
$\tau^*(k)$; to match the implementation's column names, the Monte Carlo
averages $\widehat B_\alpha$, $\widehat P_\alpha$, $\widehat D_\alpha$ are
denoted $\mathrm{AR}_\alpha$ (alarm rate), $\mathrm{PAR}_\alpha$ (premature
alarm rate), and $\mathrm{DR}_\alpha$ (detection rate), reported alongside
the same five delay statistics. $R=2500$ replications are run per
combination, exactly as in Section~\ref{sec:mc}.

\subsection{Experiment 1: common-factor volatility break}
\label{subsec:exp1}

Every series loads on factor 4 (the long-memory $\mathrm{ARFIMA}(0,0.45,0)$
factor of Table~\ref{tab:factors}) with the same fixed loading as
\eqref{eq:factorload}, $\gamma=0.5$. No level shift is injected; instead,
the i.i.d.\ $N(0,1)$ innovations that drive factor 4 through fractional
differencing, \eqref{eq:fracweights}--\eqref{eq:truncconv}, switch standard
deviation from $1$ to $\sqrt{3}$ at $\tau^*(k)$:
\begin{equation}
\varepsilon_{4,t} \sim N(0,1)\ \text{for } t < \tau^*(k), \qquad
\varepsilon_{4,t} \sim N(0,3)\ \text{for } t \ge \tau^*(k).
\label{eq:exp1vol}
\end{equation}
The idiosyncratic panel (Tables~\ref{tab:N10}--\ref{tab:N30}) is unaffected;
only the variance of the common shock changes, so the break reaches every
series solely through the factor loading.

\subsection{Experiment 2: common-factor loading break}
\label{subsec:exp2}

Each factor $j\in\{1,2,3,4\}$ is now considered separately (no volatility
break: the factor itself is generated exactly as in Table~\ref{tab:factors}
throughout). Instead, every series' loading on $f_{j,t}$ jumps from
$\gamma_{\text{pre}}=0.1$ to $\gamma_{\text{post}}=0.8$ at $\tau^*(k)$:
\begin{equation}
Y^*_{i,t} = Y_{i,t} + \gamma(t)\, f_{j,t}, \qquad
\gamma(t) = \begin{cases} 0.1, & t < \tau^*(k) \\ 0.8, & t \ge \tau^*(k), \end{cases}
\qquad i=1,\dots,N.
\label{eq:exp2load}
\end{equation}
Because each of the four factors is run as a separate combination, the
parameter grid additionally ranges over $\mathrm{factor}\in\{1,2,3,4\}$, as
in Scenario B of Section~\ref{sec:mc}.

\subsection{Experiments 3.1--3.2: heterogeneous (subgroup) factor structure}
\label{subsec:exp3}

Experiments~\ref{subsec:exp1}--\ref{subsec:exp2} load every series on the
same single factor. Experiments 3.1 and 3.2 instead partition the $N$
series into three disjoint subgroups $G_2, G_3, G_4$, each loading on a
\emph{different} factor ($f_2$, $f_3$, $f_4$ respectively) and on no other
factor (Table~\ref{tab:exp3groups}). Because no factor is shared by the
whole panel, the panel-wide pooled signal is, by construction, only
partially informative about a break confined to a single subgroup --- this
is the sense in which the factor structure is "weak" relative to
Experiments~\ref{subsec:exp1}--\ref{subsec:exp2}.

\begin{table}[h]
\centering
\caption{Subgroup composition for Experiments 3.1--3.2.}
\label{tab:exp3groups}
\begin{tabular}{cccc}
\toprule
& Factor & $N=10$ & $N=30$ \\
\midrule
$G_2$ & $f_2$ & $\{s_1,s_2,s_3\}$ & $\{s_1,\dots,s_3\}$ \\
$G_3$ & $f_3$ & $\{s_4,s_5,s_6\}$ & $\{s_4,\dots,s_{12}\}$ \\
$G_4$ & $f_4$ & $\{s_7,\dots,s_{10}\}$ & $\{s_{13},\dots,s_{30}\}$ \\
\bottomrule
\end{tabular}
\end{table}

For every replication, AR filtering is applied to the whole panel exactly
as in Section~\ref{sec:arfilter} --- each series keeps its own fixed AR$(p)$
coefficients regardless of subgroup membership. Two CUSQ tests are then
computed from the \emph{same} AR residuals: a full-panel test using
\eqref{eq:ebarclosed} with all $N$ columns, and, for each $g\in\{2,3,4\}$, a
subgroup test using \eqref{eq:ebarclosed} restricted to the $|G_g|$ columns
in $G_g$ only. The closed form of Proposition~\ref{prop:ebar} applies
verbatim to any subset of columns, so the subgroup tests require no
additional derivation; only the cross-sectional aggregation step changes,
never the AR filtering. This lets each replication compare whole-panel
monitoring against subgroup-specific monitoring for a break confined to one
subgroup.

\textbf{Experiment 3.1 (volatility break).} Each subgroup's own factor
receives the same volatility break as \eqref{eq:exp1vol} (std.\ $1\to\sqrt3$
at $\tau^*(k)$), with a fixed loading $\gamma=0.5$, applied only to the
columns in that subgroup.

\textbf{Experiment 3.2 (loading break).} Each subgroup's own factor has
constant volatility, but the subgroup's loading on it jumps from a
subgroup-specific $\gamma_{\text{pre}}$ to $\gamma_{\text{post}}$ at
$\tau^*(k)$ (Table~\ref{tab:exp32loadings}), analogous to
\eqref{eq:exp2load} but applied subgroup-by-subgroup with its own factor.

\begin{table}[h]
\centering
\caption{Pre- and post-break loadings for Experiment 3.2.}
\label{tab:exp32loadings}
\begin{tabular}{cccc}
\toprule
Subgroup & Factor & $\gamma_{\text{pre}}$ & $\gamma_{\text{post}}$ \\
\midrule
$G_2$ & $f_2$ & $0.10$ & $0.60$ \\
$G_3$ & $f_3$ & $0.25$ & $0.85$ \\
$G_4$ & $f_4$ & $0.30$ & $0.90$ \\
\bottomrule
\end{tabular}
\end{table}

For both 3.1 and 3.2, the full-panel metrics are reported with prefix
``full\_'' and each subgroup $g$'s metrics with prefix ``g$g$\_''. In
addition, for every pair of subgroups $(a,b)\in\{(2,3),(3,4),(2,4)\}$ the
gap in $5\%$-level detection rate is reported,
\begin{equation}
\Delta\mathrm{DR}_{g_a,g_b} \equiv \mathrm{DR}_{g_a,5\%} - \mathrm{DR}_{g_b,5\%},
\label{eq:deltaDR}
\end{equation}
summarizing whether a break carried by a more persistent (higher-$d$ or
higher-$\phi$) factor is detected more or less reliably than one carried by
a less persistent factor.

\subsection{Design summary}

Table~\ref{tab:targeteddesign} lists the four experiments' combinations.
$M$, factor, and the $N=30$ random-coefficient redraw rule are exactly as
in Table~\ref{tab:design}.

\begin{table}[h]
\centering
\caption{Targeted power experiment design.}
\label{tab:targeteddesign}
\begin{tabular}{lcccc}
\toprule
Experiment & $N$ & $M$ & $k$ & Factor \\
\midrule
1 (factor-4 vol.\ break)        & $\{10,30\}$ & $\{50,80,100,200,300\}$ & $\{0.2,0.3,0.5\}$ & --- \\
2 (loading break)               & $\{10,30\}$ & $\{50,80,100,200,300\}$ & $\{0.2,0.3,0.5\}$ & $\{1,2,3,4\}$ \\
3.1 (subgroup vol.\ break)      & $\{10,30\}$ & $\{50,80,100,200,300\}$ & $\{0.2,0.3,0.5\}$ & --- \\
3.2 (subgroup loading break)    & $\{10,30\}$ & $\{50,80,100,200,300\}$ & $\{0.2,0.3,0.5\}$ & --- \\
\midrule
\multicolumn{5}{l}{Combinations: $30+120+30+30=210$; $R=2500$ each; total replications $=5.25\times10^5$.} \\
\bottomrule
\end{tabular}
\end{table}

\section{Computational Details}
\label{sec:computation}

The full design requires $1.9\times10^6$ replications. Two changes to the
implementation reduce computation time substantially without changing any
statistic.

\subsection{FFT-based fractional differencing}
\label{subsec:fftfrac}

An ARFIMA process requires applying the fractional integration operator
$(1-L)^{-d}$, which has the moving-average representation
\begin{equation}
(1-L)^{-d} = \sum_{j=0}^{\infty} w_j L^j, \qquad
w_j = \frac{\Gamma(j+d)}{\Gamma(j+1)\,\Gamma(d)}.
\label{eq:fracweights}
\end{equation}
Given a short-memory series $Y_t$, the fractionally integrated series is
\begin{equation}
X_t = \sum_{j=0}^{t} w_j Y_{t-j}, \qquad t=0,1,\dots,T-1.
\label{eq:truncconv}
\end{equation}
Computing \eqref{eq:truncconv} by a direct loop costs $O(T^2)$ operations.
The same result can be obtained by the discrete convolution
\begin{equation}
(w*Y) = \mathcal{F}^{-1}\bigl[\mathcal{F}(w) \odot \mathcal{F}(Y)\bigr],
\label{eq:convthm}
\end{equation}
where $\mathcal{F}$ is the DFT computed by the FFT algorithm, at cost
$O(T\log T)$. The first $T$ elements of $(w*Y)$ equal \eqref{eq:truncconv}
exactly. For numerical stability, the weights are computed in log space:
\begin{equation}
w_j = \exp\bigl[\ln\Gamma(j+d) - \ln\Gamma(j+1) - \ln\Gamma(d)\bigr],
\label{eq:gammaln}
\end{equation}
which avoids overflow for large $j$. For $T=4000$, this reduces the
fractional differencing step from hundreds of milliseconds to under one
millisecond per series.

\subsection{Vectorized AR residuals}
\label{subsec:vectorAR}

Equation \eqref{eq:residual} computes one residual per time period. A
direct implementation loops over $t$ and computes a dot product at each
step, incurring $T-p$ separate Python function call overheads. The same
result follows from a single matrix multiplication.

\begin{lemma}
\label{lem:vector}
Let $t_0 = \max(8,\,p+1)$ and $\mathcal{I} = \{t_0,\dots,T\}$. Define
the lag matrix $Z \in \mathbb{R}^{|\mathcal{I}|\times p}$ by
\begin{equation}
Z_{t,\ell} = Y_{i,\,t-\ell}, \qquad \ell=1,\dots,p,\quad t\in\mathcal{I},
\label{eq:lagmatrix}
\end{equation}
and let $y = (Y_{i,t})_{t\in\mathcal{I}}$. Then
\begin{equation}
(\hat e_{i,t})_{t\in\mathcal{I}} = y - Z\hat\beta_{AR}.
\label{eq:vecresid}
\end{equation}
\end{lemma}

\begin{proof}
Row $t$ of $Z\hat\beta_{AR}$ equals
$\sum_{\ell=1}^p \hat\beta_{AR,\ell}\,Y_{i,t-\ell}$, which is the
predicted value in \eqref{eq:residual}. So
$(y - Z\hat\beta_{AR})_t = \hat e_{i,t}$ for every $t\in\mathcal{I}$.
\end{proof}

The matrix $Z$ is formed once and the residuals for all periods are
obtained in a single BLAS call. For $T\approx3000$ and $p\le6$, this is
about 30 times faster than the loop version, with identical numerical
output.

\subsection{Parallelisation}
\label{subsec:aggregate}

After both improvements, the time for one replication at $N=30$, $M=300$
falls from roughly $0.14$ seconds to roughly $0.02$ seconds, a reduction
of about $8\times$. Across $1.9\times10^6$ replications, this brings the
projected single-thread runtime from about 72 hours to about 9.5 hours.

The 760 parameter combinations in Table~\ref{tab:design} are independent
of one another, so each combination is assigned to one of 8 worker
processes. Each worker is seeded from OS entropy to ensure independent
random streams across processes. With 8-way parallelism, the expected
wall-clock time is 1--1.5 hours.

The unpooled pairwise aggregation of Section~\ref{sec:pool} adds an
$O(N^2M)$ Gram-matrix step per replication that was not present under
pooling. For the grid used here ($N\le30$, $M\le300$) this step is small
relative to data generation and AR filtering, so the per-replication
timings above are not expected to change materially, but they should be
re-benchmarked after the implementation update rather than assumed.

\section{Summary}
\label{sec:summary}

Scenario A produces 240 rows (40 combinations $\times$ 6 values of $k$),
each reporting $\widehat{\mathrm{Size}}_\alpha$ for
$\alpha\in\{10\%,5\%,1\%\}$. Scenario B produces 720 rows, each reporting
$\widehat{B}_\alpha$, $\widehat{P}_\alpha$, $\widehat{D}_\alpha$, and the
five delay statistics for each $\alpha$.

In every replication, the CUSQ statistic uses $D_n = C_n/C_m - (n-2)/(m-2)$
from \eqref{eq:Dn} and the boundary $g(n/m)$ from \eqref{eq:boundary},
both as specified in \citet{WangHsiao2013}.
Propositions~\ref{prop:ab} and \ref{prop:ebar} ensure that the (unpooled)
cross-pair aggregate $\bar e_t$, with every pair $(i,j)$ contributing its
own fitted slope $\hat\beta_{ij}$, is nonetheless computed exactly, at cost
$O(N^2M)+O(NT)$ instead of $O(N^2M)+O(N^2T)$ for a direct pair-by-pair
evaluation. Lemma~\ref{lem:vector} and the FFT-based
fractional-differencing step ensure that all residuals are numerically
identical to a direct implementation.

\begin{thebibliography}{9}

\bibitem[Brown et al.(1975)]{BrownDurbinEvans1975}
Brown, R.~L., J.~Durbin, and J.~M.~Evans (1975).
\newblock Techniques for testing the constancy of regression relationships
  over time.
\newblock \emph{Journal of the Royal Statistical Society, Series B}, 37,
  149--163.

\bibitem[Chu et al.(1996)]{ChuStinchcombeWhite1996}
Chu, C.-S.~J., M.~Stinchcombe, and H.~White (1996).
\newblock Monitoring structural change.
\newblock \emph{Econometrica}, 64(5), 1045--1065.

\bibitem[Inclan and Tiao(1994)]{InclanTiao1994}
Inclan, C.\ and G.~C.~Tiao (1994).
\newblock Use of cumulative sums of squares for retrospective detection of
  changes of variance.
\newblock \emph{Journal of the American Statistical Association}, 89(427),
  913--923.

\bibitem[Poskitt(2007)]{Poskitt2007}
Poskitt, D.~S.\ (2007).
\newblock Autoregressive approximation in nonstandard situations: the
  fractionally integrated and non-invertible cases.
\newblock \emph{Annals of the Institute of Statistical Mathematics}, 59,
  697--725.

\bibitem[Wang and Hsiao(2013)]{WangHsiao2013}
Wang, C.~S.-H.\ and C.~Hsiao (2013).
\newblock Real-time monitoring test for realized volatility.
\newblock \emph{Journal of Time Series Econometrics}, 5(1), 1--24.

\end{thebibliography}

\end{document}
